% ============================================================
% 分片：s02a  对应原文 PDF 第 3--6 页
% 负责章节：§2 Twisting & mirror symmetry；2.1 A modern introduction to twisting；2.1.1 The nilpotence variety
% 译者：translator-01   审校状态：初稿
% ============================================================

% ==== 原文 PDF 第 3 页 ====
\section{扭变与镜像对称}

\begin{keytakeaway}
\textbf{核心物理图像：}扭变是在超对称理论中选取一个幂零超荷 $Q$，只保留 $Q$ 闭的局域算子并对 $Q$ 恰当算子作等价，
从而把理论约化到由 $Q$ 上同调所刻画的 BPS 可观测量；当应力张量均为 $Q$ 恰当时，关联函数与度规无关，得到拓扑（或全纯–拓扑）理论。

\textbf{数学工具：}取 $Q$ 上同调（$Q$ 闭模去 $Q$ 恰当）；幂零簇——可容许 $Q$ 构成的射影簇——及其按 $\Gamma$ 矩阵秩的分层；
$\mathrm{Spin}(d)$ 与 $R$ 对称群对幂零簇的作用。

\textbf{主要结论：}二维 $\mathcal{N}=(2,2)$ 的幂零簇是四个 $\mathbb{P}^1$；秩 1 分支给出全纯的「半扭变」，
秩 2 分支给出拓扑的 $A$ 扭变与 $B$ 扭变。
\end{keytakeaway}

尽管对偶并不必然源于超对称的存在（确实，许多非常有趣的对偶并非如此！回忆 Kramers--Wannier 对偶；
关于更多现代的例子，%
% ==== 原文 PDF 第 4 页 ====
% [跨页接续: 原文第3页/4页，该句自上一页续接]
可参见 [6--8] 等），但由于 BPS 量无处不在、且我们可以在两个候选的对偶框架中分别计算它们，
超对称是启发并检验对偶的特别方便的方式。
借助一种称为\emph{扭变}（twisting）的程序，可以便捷地分离出这些完整物理对偶中的受保护部分。
扭变把数学上尚未良定义的量局域化为定义良好的对象，而这些对象仍然保留了大量的物理与数学内涵。
我们首先介绍扭变，它既适用于超对称场论，也适用于超弦理论。
我们将特别关注镜像对称这一与数学联系紧密的弦对偶——由于扭变程序所提供的精确性，对它的研究已蓬勃发展。
更多细节可参阅两部详尽的教材 [9, 10]，我们在若干地方将遵循它们的表述。

\subsection{扭变的现代导引}

扭变是一种可以在超对称场论上施行的操作，粗略地说，它把所考虑的物理可观测量的子空间限制到一个更简单、更易处理的（BPS）子集。
本节的讨论非常一般地适用于超对称场论，不过我们很快就会有兴趣把它应用于描述某些超弦世界面的二维超对称场论。
对具有引力动力学的时空理论直接施行扭变要更微妙，是一个超出本讲义范围的活跃研究方向，但可参见例如 [11--14]。
我们主要在欧几里得号差下工作，在那里扭变得到最好的理解且定义良好。

扭变始于选取一个幂零超荷，即所关心的超对称代数中平方为零 $[Q,Q]=0$ 的一个超荷 $Q$。
\footnote{我们采用这样的约定：$[-, -]$ 是分次对易子，即 $[a,b]=ab-(-1)^{F(a)F(b)}ba$；
若 $a$ 或 $b$ 之一为玻色子（$F(a)$ 与 $F(b)$ 中至少一个为偶数），它是通常的对易子；
否则（$F(a)$ 与 $F(b)$ 均为奇数）它是反对易子。
它对费米子宇称是分次反对称的，$[a,b]=-(-1)^{F(a)F(b)}[b,a]$，并满足分次版本的 Jacobi 恒等式
$[a,[b,c]]=[[a,b],c]+(-1)^{F(a)F(b)}[b,[a,c]]$；
该式表明线性映射 $[a,-]$ 是括号 $[-, -]$ 的一个导子（与 $a$ 具有相同的统计性质）。}
这是一种与 BRST 对称性非常相似的费米型对称性，其行为方式也相同。
我们考虑 $Q$ 不变量，于是经过扭变后存活的算子是 $Q$ 闭的：$[Q,O]=0$。
若进一步假设超对称 $Q$ 没有自发破缺，即 $Q$ 湮灭真空，则可推出 $Q$ 闭的 $O$ 与其他 $Q$ 闭局域算子的关联函数在 $O\to O+[Q,\Lambda]$ 下不变：
\begin{equation}\tag{2.1}
\langle (O+[Q,\Lambda])\dots\rangle=\langle O\dots\rangle+\langle [Q,\Lambda\dots]\rangle=\langle O\dots\rangle.
\end{equation}
% [Tutor Note: 推导补足：$\langle [Q,X]\rangle=\langle 0|[Q,X]|0\rangle=\langle 0|QX|0\rangle\pm\langle 0|XQ|0\rangle=0$，
% 其中用到超对称不破缺的假设 $Q|0\rangle=0$、$\langle 0|Q=0$。这就是「$Q$ 恰当算子的关联函数为零」的来源，
% 也是下文 (2.2)、(2.3) 中一切「$Q$ 恰当故为零」论证的依据。]
这蕴含着等价关系 $O\sim O+[Q,\Lambda]$，其中形如 $[Q,\Lambda]$ 的算子称为 $Q$ 恰当的（$Q$-exact）。
因此，只要一个量子场论是由其关联函数的数据来定义的，扭变理论中的局域算子就是 $Q$ 闭的局域算子，
模去 $Q$ 恰当局域算子的添加：它们由 $Q$ 上同调的元素标记。

设想我们以 BRST 形式量子化一个具有某种规范对称性的理论，从而关于某个 BRST 微分 $d$ 取上同调。
用这套语言表述扭变的一个简单办法是：把 BRST 微分 $d$ 用这个新超荷 $Q$ 增广，并关于形变后的微分 $d_Q=d+Q$ 取上同调：
扭变理论中的算子由 $d_Q$ 闭的算子给出，模去 $d_Q$ 恰当的算子。
作一初步近似，扭变理论中的算子可以看成是规范/BRST 不变的%
% ==== 原文 PDF 第 5 页 ====
% [跨页接续: 原文第4页/5页，该句自上一页续接]
（$d$ 闭）算子，同时又是 $Q$ 不变的（$Q$ 闭）。
这并不完全正确：必须修正上述答案，以计入那些既非 $d$ 闭、也非 $Q$ 闭、却仍是 $d_Q$ 闭的算子。
% [原文疑误: 原文 "can be though of as" 应为 "can be thought of as"（漏写字母 t）；语义明确，译文按 "thought" 处理。]
% [Tutor Note: 推导补足：$d_Q^2=\{d,Q\}$，故只有在 $\{d,Q\}=0$（$d$ 与 $Q$ 反对易）时 $d_Q$ 才幂零、其上同调才有定义；
% 原文此处未明言这一前提。]
有一种以同调方式计算这些修正的办法，即在每一步计算一个适当的上同调，这称为\emph{谱序列}（spectral sequence）。
% [Tutor Note: 谱序列的机制：按 $d$ 过滤 $d_Q$，先取 $E_1=H_d$，$Q$ 在 $E_1$ 上诱导微分 $d_1$，逐次取上同调得 $E_r\Rightarrow H_{d_Q}$。
% 若 $\{d,Q\}=0$，$E_2$ 项大致由 $d$ 与 $Q$ 的混合上同调给出；原文提醒：$H_d\cap H_Q$ 只是完整答案的近似，扭变理论一般并不对应原理论中简单的可观测量子集。]
计算 $d$ 上同调与 $Q$ 上同调的交集可以是完整答案的一个有用近似，
但一般而言我们提醒读者：扭变理论并不总能捕捉到原理论中一个简单的可观测量子集。

\subsubsection{幂零簇}

对给定的超对称代数，幂零超荷 $Q$ 往往有很多个，因而可能的扭变也有很多。
设共有 $N$ 个超荷 $Q_i$，$i=1,\dots,N$（$i$ 是旋量与 $R$ 对称指标的某种组合），其括号为 $[Q_i,Q_j]=\Gamma^\mu_{ij}P_\mu$，
其中 $P_\mu$ 是时空动量算子；若 $Q=q^iQ_i$，则 $Q$ 的幂零性，即 $[Q,Q]=0$，转化为关于 $q^i$ 的一组二次方程：
$\Gamma^\mu_{ij}q^iq^j=0$，$\mu=1,\dots,d$。
我们要求至少有一个 $q^i$ 非零（否则 $Q=0$），并注意到若 $Q$ 满足此方程，则对任意 $\lambda\in\mathbb{C}^*$，$\lambda Q$ 也满足；
于是，给定超对称代数中幂零超荷的模空间——有时称为该代数的\emph{幂零簇}（nilpotence variety）——
自然地是 $(N-1)$ 维射影空间 $\mathbb{P}^{N-1}$ 的一个闭子簇 [15, 16]。
% [背景补充: 幂零簇是「可容许扭变」的模空间：齐次坐标 $q^i$ 整体标度等价（同一个 $Q$ 只差一个非零常数倍），
% 故取射影空间；固定超代数后，分类其幂零簇就等于分类该超代数所有可能的扭变。]

幂零簇通常是奇异的，但它带有 $\mathrm{Spin}(d)$ 的复化与超对称代数的 $R$ 对称群 $G_R$ 的复化的自然作用，
后者转动齐次坐标 $q^i$ 的指标 $i$。
此外，它还具有由 $d\times N$ 矩阵 $\Gamma^\mu_{ij}q^i$ 的秩给出的自然分层，
即按属于映射 $[Q,-]$ 的像的平移 $P_\mu$ 的数目分层。
落在 $[Q,-]$ 的像中的平移 $P_\mu$ 按定义是 $Q$ 恰当的，因而在扭变理论中必然平凡：
若存在算子 $Q_\mu$ 使得 $[Q,Q_\mu]=iP_\mu\sim\partial_\mu$，则 $Q$ 闭可观测量的关联函数沿 $x^\mu$ 方向局部为常数。
显式地（着眼于 $x$ 处的单个插入）：
\begin{equation}\tag{2.2}
\begin{aligned}
\partial_\mu\langle O(x)\dots\rangle &\sim \langle [Q,[Q_\mu,O(x)]]\dots\rangle\\
&\sim \langle [Q,[Q_\mu,O(x)]\dots]\rangle=0
\end{aligned}
\end{equation}
% [Tutor Note: 推导补足：由 $[Q,Q_\mu]=iP_\mu$ 得 $P_\mu=-i[Q,Q_\mu]$；
% 对 $Q$ 闭的 $O$，$\partial_\mu O=[P_\mu,O]=-i[[Q,Q_\mu],O]$，再用分次 Jacobi 恒等式并把 $[Q,O]=0$ 的项去掉，
% 即得 $\partial_\mu O\sim[Q,[Q_\mu,O]]$。于是 $\partial_\mu\langle O\dots\rangle$ 是 $Q$ 恰当算子的关联函数，由 (2.1) 的论证为零。]
% [Tutor Note: 物理图像：分层按被平凡化的平移数目 $r$ 划分——$r=0$ 时没有平移被平凡化；$r=d$ 时所有平移都 $Q$ 恰当，
% 得到完全拓扑扭变；中间情形 $0<r<d$ 即全纯或混合扭变，扭变后只在部分方向上具有拓扑不变性。]
算子 $Q_\mu$ 将在第 3 节再次出现，并支撑扭变量子场论中可观测量的大部分同伦代数结构。

各种超对称代数中 $\Gamma$ 矩阵的结构意味着，对于给定的 $Q$，被平凡化的那些平移会分成两类：
适当选取坐标后会发现，例如对 $\mu=1,\dots,n\le d$，$P_\mu$ 被平凡化，而只有复组合
$P_{\bar a}=\frac{1}{2}(P_{n+2a-1}+iP_{n+2a})$（$a=1,\dots,m$）被平凡化，其中 $d=n+2m$。
由此，$Q$ 的上同调在沿 $x^\mu$ 的任意平移下不变，在沿反全纯平移
$\bar z^{\bar a}:=x^{n+2a-1}-ix^{n+2a}$ 下也不变。
最简单且研究最多的幂零超荷 $Q$ 是 $n=d$ 的情形，它们称为\emph{拓扑超荷}，相应的扭变称为\emph{拓扑扭变} [17]。
更一般地，$Q$ 称为\emph{全纯–拓扑的}或\emph{混合超荷}；当 $d$ 为偶数、$d=2m$ 时，$Q$ 称为\emph{全纯超荷} [18, 19]，
相应的扭变也依此命名。
% [原文疑误: 原文 "these are called a topological supercharges" 中冠词与复数不一致，应为 "these are called topological supercharges"；译文按复数处理。]

% ==== 原文 PDF 第 6 页 ====
拓扑扭变的理论往往具有如下性质：应力张量的所有分量（它们支配度规的变分）都是 $Q$ 恰当的：
\begin{equation}\tag{2.3} T_{\mu\nu}=\frac{\delta S}{\delta h^{\mu\nu}}=[Q,G_{\mu\nu}]. \end{equation}
% [Tutor Note: 由 $\delta S=\int T_{\mu\nu}\,\delta h^{\mu\nu}$（原文约定）得 $\delta\langle O_1\cdots O_n\rangle=\langle \delta S\,O_1\cdots O_n\rangle$；
% 把 $T_{\mu\nu}=[Q,G_{\mu\nu}]$ 代入，再用 (2.1) 的「$Q$ 恰当插入为零」即得度规无关性。]
% [Tutor Note: 易错预警：通常的应力张量定义是 $T^{\mu\nu}=\frac{2}{\sqrt{g}}\frac{\delta S}{\delta h_{\mu\nu}}$，与原文 (2.3) 的约定相差 $2/\sqrt{g}$ 因子且指标位置相反；
% 与 (2.3) 配套使用时应统一为对 $h^{\mu\nu}$ 变分，否则会漏掉因子。]
在实践中实现这一条件的常见方式是：整个拉格朗日量本身就可以写成 $Q$ 的对易子。
由这一条件，理论中关联函数的度规无关性，即 $\dfrac{\delta\langle O_1\dots O_n\rangle}{\delta h^{\mu\nu}}=0$，是直接的，
因为作用量的变分会带下一个 $T_{\mu\nu}$ 因子，而它是 $Q$ 恰当的，因而在关联函数中为零。
混合的全纯–拓扑理论具有类似的性质。
从启发式看，人们预期在拓扑方向上得到拓扑场论（TQFT），在全纯方向上得到全纯量子场论；
后者在二维情形具有顶点代数或手征代数的结构（这在共形场论中已为人熟知），而在更高维情形则由手征代数的「高阶」类比给出。
我们将在后文进一步阐述高阶结构。

作为贯穿全篇的例子，考虑二维 $\mathcal{N}=(2,2)$ 超对称代数。
我们在欧几里得号差下工作，并在世界面上取复坐标 $z,\bar z$。
超对称代数由 $N=4$ 个超荷 $Q_\pm,\bar Q_\pm$ 生成，其非零括号（无中心荷时）为
\begin{equation}\tag{2.4}
[Q_+,\bar Q_+]=2P_{\bar z}\qquad [Q_-,\bar Q_-]=-2P_z
\end{equation}
其中 $P_z,P_{\bar z}$ 为全纯/反全纯动量。
\footnote{在 Lorentz 号差下，我们有左/右移动量 $P_\pm$，其反对易子为 $[Q_\pm,\bar Q_\pm]=2P_\pm$。
Lorentz 号差下的左/右移坐标 $x^\pm=x^0\pm x^1$ 经 Wick 转动后，与欧几里得全纯坐标的关系为 $x^+=\bar z$ 及 $x^-=-z$。}
该超代数在 $U(1)_V$ 矢量 $R$ 对称转动（由荷 $R_V$ 生成）下不变，后者对超荷的作用为
\begin{equation}\tag{2.5} [R_V,Q_\pm]=Q_\pm,\qquad [R_V,\bar Q_\pm]=-\bar Q_\pm, \end{equation}
并且在 $U(1)_A$ 轴 $R$ 对称转动（由 $R_A$ 生成）下也不变，其作用为
\begin{equation}\tag{2.6} [R_A,Q_\pm]=\pm Q_\pm,\qquad [R_A,\bar Q_\pm]=\mp\bar Q_\pm. \end{equation}
% [Tutor Note: 易错预警：$R_V$ 给 $Q_\pm$ 与 $\bar Q_\pm$ 相反的荷（(2.5)），$R_A$ 给 $Q_\pm$ 以 $\pm$、给 $\bar Q_\pm$ 以 $\mp$（(2.6)）；
% 这些符号是约定，须与 (2.7) 的幂零条件及 $A/B$ 扭变的定义一并核对；换用其他文献的约定时符号可能相反。]
超荷 $Q=q^\alpha Q_\alpha+\bar q^\alpha\bar Q_\alpha$ 幂零，当且仅当
\begin{equation}\tag{2.7}
q^+\bar q^+=0\qquad q^-\bar q^-=0.
\end{equation}
因此，该超对称代数的幂零簇可以等同于四个在极点处两两相切的 $\mathbb{P}^1$；见 Figure 1。
% [Tutor Note: (2.7) 的两式各自是「二选一」的零条件：对指标 $+$ 的一对 $(q^+,\bar q^+)$ 二选一，对指标 $-$ 的一对 $(q^-,\bar q^-)$ 二选一，
% 故共 $2\times2=4$ 个分支；每个分支上两个自由分量之比给出 $\mathbb{P}^1$ 的齐次坐标，四个 $\mathbb{P}^1$ 在四个坐标点（各分支的「极点」）处两两相切。]
不计代数的对称性，例如整体标度、宇称 $+\leftrightarrow-$，或电荷共轭 $Q_\pm\leftrightarrow\bar Q_\pm$，我们总可以取 $\bar q^+=1$（从而 $q^+=0$）。
于是有三种可能：1）$\bar q^-=0=q^-$；2）$\bar q^-=0$ 而 $q^-\neq0$；3）$\bar q^-\neq0$ 而 $q^-=0$。
容易看出，第一种情形 $Q=Q_H=\bar Q_+$ 是一个秩 1 的全纯超荷。
相应的扭变历史上称为\emph{半扭变}（half twist）[20--22]。
% [Tutor Note: 物理图像：$Q_H=\bar Q_+$ 只使一半的超荷成为标量（另一半仍为旋量），故称「半扭变」；所得理论只在一个复方向上拓扑化，
% 因此是全纯理论而非完全拓扑理论。$Q_A,Q_B$ 则是秩 2 的拓扑超荷，给出完整的 $A/B$ 拓扑扭变。]
对于第二种（相应地第三种）情形，我们总可以借助（复化的）$R$ 对称转动与重新标度取 $q^-=1$（相应地 $\bar q^-=1$），
从而看出本质上只有一种选择 $Q_A=\bar Q_++Q_-$（相应地 $Q_B=\bar Q_++\bar Q_-$），且它是秩 2 的拓扑超荷。
所得的拓扑扭变相应地称为 $A$ 扭变与 $B$ 扭变。
